% Copyright (c) 2004-2010 Marc Fonvieille
% All rights reserved.
%
% Redistribution and use in source and binary forms, with or without
% modification, are permitted provided that the following conditions
% are met:
% 1. Redistributions of source code must retain the above copyright
%    notice, this list of conditions and the following disclaimer.
% 2. Redistributions in binary form must reproduce the above copyright
%    notice, this list of conditions and the following disclaimer in the
%    documentation and/or other materials provided with the distribution.
%
% THIS SOFTWARE IS PROVIDED BY THE AUTHOR AND CONTRIBUTORS ``AS IS'' AND
% ANY EXPRESS OR IMPLIED WARRANTIES, INCLUDING, BUT NOT LIMITED TO, THE
% IMPLIED WARRANTIES OF MERCHANTABILITY AND FITNESS FOR A PARTICULAR PURPOSE
% ARE DISCLAIMED.  IN NO EVENT SHALL THE AUTHOR OR CONTRIBUTORS BE LIABLE
% FOR ANY DIRECT, INDIRECT, INCIDENTAL, SPECIAL, EXEMPLARY, OR CONSEQUENTIAL
% DAMAGES (INCLUDING, BUT NOT LIMITED TO, PROCUREMENT OF SUBSTITUTE GOODS
% OR SERVICES; LOSS OF USE, DATA, OR PROFITS; OR BUSINESS INTERRUPTION)
% HOWEVER CAUSED AND ON ANY THEORY OF LIABILITY, WHETHER IN CONTRACT, STRICT
% LIABILITY, OR TORT (INCLUDING NEGLIGENCE OR OTHERWISE) ARISING IN ANY WAY
% OUT OF THE USE OF THIS SOFTWARE, EVEN IF ADVISED OF THE POSSIBILITY OF
% SUCH DAMAGE.
%
% $FreeBSD: doc/de_DE.ISO8859-1/flyer/flyer.tex,v 1.8 2010/01/24 18:49:24 bcr Exp $
% $FreeBSDde: de-docproj/flyer/flyer.tex,v 1.10 2010/01/23 20:27:46 bcr Exp $
% basiert auf: 1.17
%
% FreeBSD Flyer
% Use make FORMAT (with FORMAT: pdf, ps or dvi) to build the flyer.
% Two layouts are available: one using Beastie the other one using the
% FreeBSD Logo.  The layout selection is done below via the value of
% the \logo variable.  By default the Beastie layout is enabled.
%
\documentclass[11pt]{article}
\usepackage[T1]{fontenc}
\usepackage[latin1]{inputenc}
% Use the right language
\usepackage[german]{babel}
\usepackage{pslatex}
\usepackage{graphicx}
\usepackage{fancybox}
\usepackage{url}
% Use the right papersize, do not forget to change also the Makefile
\usepackage[verbose,a4paper,noheadfoot,margin=1cm]{geometry}

\usepackage{ifthen}
% Use of the Logo (set the \logo variable below to true) or Beastie
% (\logo variable set to false).
\newcommand{\logo}{false}

% Colors settings
\usepackage{color}
\ifthenelse{\equal{\logo}{true}}{
  \definecolor{bkgrdtitle}{rgb}{.69,0,0}
  \definecolor{redtitle}{rgb}{.65,.16,.22}
  \definecolor{ovalboxcolor}{rgb}{.69,0,0}
}
{
\definecolor{bkgrdtitle}{rgb}{1,.84,.22}
\definecolor{redtitle}{rgb}{.65,.16,.22}
\definecolor{ovalboxcolor}{rgb}{.65,.16,.22}
}

% Some macros
\ifthenelse{\equal{\logo}{true}}{
  \newcommand{\titledframe}[3]{%
  \boxput*(0,1){\colorbox{bkgrdtitle}{\color{white} \large{\textbf{\textsf{#1}}}}} {\setlength {\fboxsep}{12pt} \color{ovalboxcolor}\Ovalbox {\color{black}\begin{minipage}{#3}#2\end{minipage}}}
}}
{
\newcommand{\titledframe}[3]{%
\boxput*(0,1){\colorbox{bkgrdtitle}{\color{black} \large{\textbf{\textsf{#1}}}}} {\setlength {\fboxsep}{12pt} \color{ovalboxcolor}\Ovalbox {\color{black}\begin{minipage}{#3}#2\end{minipage}}}
}}

\newcommand{\reg}{$^{\mbox{\tiny \textregistered}}$}
\newcommand{\tm}{$^{\mbox{\tiny TM}}$}

\newenvironment{itemizeflyer}%
{ \begin{list}%
	{\textendash}%
	{ \setlength{\leftmargin}{5pt}%
	  \setlength{\itemsep}{0pt}%
	  \setlength{\parskip}{0pt}%
	  \setlength{\parsep}{0pt}}}
{ \end{list}}

\pagestyle{empty}

\begin{document}

\begin{center}
\ifthenelse{\equal{\logo}{true}}{
\includegraphics[scale=0.4]{logo-full.eps}
\vspace{1mm}
}
{
\fontsize{40}{36}\selectfont
{\color{redtitle} \textrm{\textbf{FreeBSD}}}\medskip}
\end{center}
%\vspace{2mm}

% Main part
\begin{center}
\ifthenelse{\equal{\logo}{true}}{
\newcommand{\size}{17.3cm}
}
{\newcommand{\size}{12.7cm}}

\titledframe{Was ist FreeBSD?}{
FreeBSD ist ein modernes Betriebssystem f�r x86-kompatible (einschlie�lich
Pentium\reg und Athlon\tm), amd64-kompatible (einschlie�lich Opteron\tm,
Athlon\tm 64 und EM64T), UltraSPARC\reg-, IA-64- (Intel\reg Itanium\reg
Prozessor-Familie), PC-98- sowie ARM-Architekturen.

FreeBSD stammt von BSD, dem UNIX\reg-Betriebssystem der University of
California, Berkeley, ab.
}{\size}
\ifthenelse{\equal{\logo}{false}}{
\begin{minipage}{4cm}
\includegraphics[scale=0.3]{../../share/images/flyer/beastie.eps}
\end{minipage}
}

\ifthenelse{\equal{\logo}{true}}{
\vspace{5mm}}{\vspace{3mm}}

\titledframe{Herausragende Funktionen}{
In den Bereichen Netzwerk, Leistungsf�higkeit, Sicherheit und
Kompatibilit�t besitzt FreeBSD heute schon Funktionen, die in
anderen Betriebssystemen, selbst in den besten kommerziellen,
fehlen.
}{5cm}
\titledframe{\textsf{\textbf{Leistungsf�hige Internet-Dienste}}}{
FreeBSD enth�lt eine von vielen Leuten als Referenzimplementierung
betrachtete TCP/IP-Software, den 4.4BSD TCP/IP-Protokoll-Stack.
Es ist daher ideal f�r Netzwerkanwendungen und das Internet.
FreeBSD ist bestens geeignet f�r Internet- oder Intranet-Server.
Auch unter h�chsten Lasten arbeiten die Netzwerkdienste zuverl�ssig.
Der effiziente Umgang mit dem Speicher garantiert schnelle Antwortzeiten
f�r tausende gleichzeitig laufende Benutzerprozesse.
}{11.12cm}

\vspace{5mm}

\titledframe{W�hlen Sie aus einer Vielzahl von Anwendungen!}{
FreeBSD wird mit mehr als 21.000 installationsfertigen Softwarepaketen
ausgeliefert, unter anderem: Apache, Samba, MySQL\reg,
OpenOffice.org, KDE, GNOME, MPlayer, etc.
\medskip

FreeBSD bietet die Ports-Sammlung, ein System, das Quellcode
aus dem Internet oder von einer CD-ROM herunterl�dt und daraus
Bin�rdateien kompiliert, inklusive dem Bau und der Installation von
abh�ngigen Paketen.  Alle Operationen sind dabei f�r den Benutzer
transparent.
\medskip

Die meisten GNU/Linux-Applikationen sind auch unter FreeBSD lauff�hig.
So laufen etwa Acrobat Reader\reg, Oracle\reg, SAP/R3\reg, Mathematica\reg,
Quake3, etc. ohne merkliche Geschwindigkeitsverluste.
}{17.3cm}

\vspace{5mm}

\titledframe{FreeBSD ist einfach zu installieren}{
FreeBSD kann von verschiedenen Medien, wie CD-ROM, DVD-ROM, Disketten,
Bandlaufwerken, MS-DOS\reg Partitionen, oder auch �ber eine
Netzwerkverbindung (per anonymem FTP oder NFS) installiert werden.
}{17.3cm}

\vspace{5mm}

\titledframe{FreeBSD ist frei}{
FreeBSD ist kostenlos verf�gbar und wird mit vollst�ndigem Quellcode
ausgeliefert.  Der Gro�teil dieses Quellcodes steht unter der
BSD-Lizenz.  Im Gegensatz zur GPL-Lizenz des Linux Kernels erlaubt es
die BSD-Lizenz, ver�nderte Versionen des Originalquellcodes als
Bin�rdateien zu verbreiten, ohne den Quellcode mitliefern zu m�ssen.
Dies f�hrt dazu, dass Firmen den FreeBSD-Quellcode als Basis f�r ihre
Produkte nutzen.  Oft flie�en dabei Teile des weiterentwickelten
Codes zur�ck an die Gemeinschaft.
}{17.3cm}

\vspace{5mm}

\titledframe{F�r nahezu jede Aufgabe geeignet}{
FreeBSD ist f�r eine gro�e Anzahl von Desktop-, Server-, Embedded- und
Applianceanwendungen gut geeignet.
Heutzutage ist FreeBSD nicht nur ein Betriebssystem f�r Server,
sondern bietet sich auch f�r Endnutzer, vor allem Umsteiger von
Windows\reg und GNU/Linux, an.
}{17.3cm}

\vspace{5mm}

\titledframe{Kontakt}{
\begin{itemizeflyer}
\item Webseite: \url{http://www.FreeBSD.org/de/}
\item FreeBSD-Handbuch: \url{http://www.FreeBSD.org/de/handbook/}
\end {itemizeflyer}
}{17.3cm}
\end{center}
%
% Copyrights
\begin{center}
\tiny \copyright 2004-2010 The FreeBSD Project\\
\ifthenelse{\equal{\logo}{true}}{
FreeBSD und das FreeBSD-Logo sind registrierte Markenzeichen
der FreeBSD Foundation.\\}
{FreeBSD ist ein registriertes Markenzeichen der FreeBSD Foundation.\\}
Alle anderen Firmen- und Produktnamen k�nnen Markenzeichen der
jeweiligen Firmen sein.\\
\ifthenelse{\equal{\logo}{false}}{
BSD Daemon, \copyright 1988 von Marshall Kirk McKusick.  Alle Rechte
vorbehalten.}{}
\end{center}

\end{document}
